% Figure: energy decay curve (Schroeder integral) of a room impulse response,
% showing the -60 dB line and the extrapolated T60 from the -5 to -35 dB slope.
\begin{figure}[htbp]
\centering
\begin{tikzpicture}
\begin{axis}[
  width=12cm, height=7.5cm,
  xlabel={time after source stops (s)}, ylabel={energy decay level (dB)},
  xmin=0, xmax=1.0, ymin=-70, ymax=2,
  axis lines=left,
  legend pos=north east,
  legend cell align=left,
  every axis x label/.style={at={(current axis.right of origin)}, anchor=west},
  every axis y label/.style={at={(current axis.above origin)}, anchor=south},
  grid=major, grid style={enggray!20},
  tick label style={font=\scriptsize},
  label style={font=\small},
  legend style={font=\scriptsize},
]
% idealised linear decay at 120 dB/s -> T60 = 0.5 s, with early-time curvature
\addplot[engblue, very thick, domain=0:0.85, samples=120]
  {-120*x + 3*sin(deg(20*x))*exp(-3*x)};
\addlegendentry{measured decay curve}
% fitted -5 to -35 dB line extrapolated to -60 dB
\addplot[engred, very thick, dashed, domain=0:0.5, samples=2] {-120*x};
\addlegendentry{fit (5 to 35 dB), $T_{60}=0.5$ s}
% the -60 dB reference line
\addplot[enggray, thin, domain=0:1, samples=2] {-60};
\node[enggray, font=\scriptsize, anchor=south west] at (axis cs:0.02,-59) {$-60$ dB};
% T60 marker
\addplot[engred, only marks, mark=*, mark size=2pt] coordinates {(0.5,-60)};
\end{axis}
\end{tikzpicture}
\caption[Energy decay and reverberation time]{Backward-integrated
energy decay curve (the Schroeder integral) of a room impulse response,
the standard route to reverberation time. The measured decay (blue) is
noisy near the floor and slightly curved at early times, so $T_{60}$ is
read not from the raw $-60\,\si{\decibel}$ crossing but from a straight
line fitted over the cleaner $-5$ to $-35\,\si{\decibel}$ range and
extrapolated (red dashed) to $-60\,\si{\decibel}$. Here the fit gives a
$120\,\si{\decibel\per\second}$ slope and $T_{60}=0.5\,\si{\second}$,
typical of a furnished living room. Fitting over a defined range rather
than chasing the noisy tail is the discipline that makes $T_{60}$
repeatable across measurements. The reader produces a plot like this in
the chapter project.}
\label{fig:vol03ch07:rt60-decay}
\end{figure}
